\documentclass[11pt]{amsart}

\usepackage[T1]{fontenc}
\usepackage{lmodern}
\usepackage{microtype}
\usepackage{amsmath,amssymb,amsthm}
\usepackage{booktabs}
\usepackage[hidelinks]{hyperref}

\hypersetup{
  pdftitle={The p=3 Case of a Conjecture of Chauhan, Shukla, and Vinayak},
  pdfsubject={Cut complexes and shellability},
  pdfkeywords={cut complex, shellability, power of a cycle, computer-assisted proof}
}

\newtheorem{theorem}{Theorem}
\newtheorem{proposition}[theorem]{Proposition}
\newtheorem{lemma}[theorem]{Lemma}
\theoremstyle{remark}
\newtheorem{remark}[theorem]{Remark}

\newcommand{\Z}{\mathbb Z}
\newcommand{\Sph}{\mathbb S}
\newcommand{\triple}[3]{\{#1,#2,#3\}}

\title[The $p=3$ case of a 3-cut conjecture]
{The $p=3$ Case of a Conjecture of
Chauhan, Shukla, and Vinayak}

\date{}

\subjclass[2020]{Primary 05E45; Secondary 05C69, 55U05}
\keywords{cut complex, shellability, power of a cycle, computer-assisted proof}

\begin{document}

\begin{abstract}
We prove the $p=3$ case of Conjecture~4.2(ii) of Chauhan, Shukla,
and Vinayak on $3$-cut complexes of powers of cycles.  Within the
conjectured $p=3$ range, the only cases not covered by the theorem of Chauhan
and Shukla are $n=13$ and $n=14$.
By locally reordering, in their explicit facet order, a block of six facets
(five of which change position) and a block of ten, we obtain shellings with
one and eight spanning facets.  Hence
\[
  \Delta_3(C_{13}^3)\simeq \Sph^9,
  \qquad
  \Delta_3(C_{14}^3)\simeq \bigvee^8 \Sph^{10}.
\]
The finite verification is exact.  Both reordered blocks are printed in full
below; the underlying facet order is the one defined by Chauhan and Shukla,
specialized to $p=3$, and is not reproduced here.
\end{abstract}

\maketitle

\section{Main result}

For a graph $G$ on vertex set $V$, the $3$-cut complex $\Delta_3(G)$ is the
pure simplicial complex whose facets are the sets $V\setminus T$, where
$T\subseteq V$ has cardinality three and $G[T]$ is disconnected
\cite{BayerEtAl,CSV}.  Write $C_n^p$ for the $p$th power of the cycle $C_n$.

Conjecture~4.2(ii) of Chauhan, Shukla, and Vinayak asserts that for $p\ge3$
and $n\ge4p+1$,
\begin{equation}\label{eq:conjecture}
  \Delta_3(C_n^p)
  \simeq
  \bigvee^{(n^2-4np-n+2)/2}\Sph^{n-4},
\end{equation}
and, more strongly, that the complex is shellable
\cite[Conjecture~4.2(ii)]{CSV}.  Chauhan and Shukla proved this for
$n\ge6p-3$ \cite[Theorem~1.4]{CS}.  At $p=3$, their result begins at
$n=15$, whereas the conjectured range begins at $n=13$.

\begin{theorem}\label{thm:main}
Part~\textup{(ii)} of Conjecture~4.2 of Chauhan--Shukla--Vinayak holds for
$p=3$.  Thus, for every $n\ge13$, the complex $\Delta_3(C_n^3)$ is
shellable and
\[
  \Delta_3(C_n^3)
  \simeq
  \bigvee^{(n^2-13n+2)/2}\Sph^{n-4}.
\]
\end{theorem}

The proof is the two explicit certificates of Section~2, for $n=13$ and
$n=14$, combined with \cite[Theorem~1.4]{CS} for $n\ge15$.  The infinite part
of the range is thus that theorem's, and \cite{CS} is at the time of writing
an unrefereed preprint; the content contributed here is exactly the two
boundary cases.  Part~\textup{(i)} of Conjecture~4.2, and the cases $p\ge4$
of part~\textup{(ii)}, are untouched.

The integral homology ranks for the two boundary cases already appear in
Table~2 of \cite{CSV}.  The new point is an explicit shelling certificate,
which proves shellability and the asserted homotopy types.

\section{The two boundary certificates}

Let $V_n=\Z/n\Z=\{0,1,\ldots,n-1\}$ and
\[
  d_n(a,b)=\min\{|a-b|,n-|a-b|\}.
\]
Two vertices are adjacent in $C_n^3$ exactly when $d_n(a,b)\le3$.
We encode a facet $F=V_n\setminus T$ by its omitted triple $T$.
A graph on three vertices is disconnected exactly when it has at most one
edge; hence the omitted triples are precisely the triples containing at most
one pair at circular distance at most three.

For a shelling $F_1,\ldots,F_m$, a facet $F_j$ is \emph{spanning} if
$\partial F_j\subseteq\bigcup_{i<j}F_i$.
The following elementary criterion is convenient for exact verification.

\begin{lemma}\label{lem:criterion}
Let $T_1,\ldots,T_m$ be distinct omitted triples, set
$F_j=V_n\setminus T_j$, and, for $j\ge2$, define
\[
  R_j=\left\{
  x\in V_n\setminus T_j:
  T_r=(T_j\setminus\{y\})\cup\{x\}
  \text{ for some }r<j\text{ and }y\in T_j
  \right\}.
\]
Then $F_1,\ldots,F_m$ is a shelling order if and only if
\begin{equation}\label{eq:exchange}
  (T_i\setminus T_j)\cap R_j\ne\varnothing
  \qquad\text{for every }i<j.
\end{equation}
Moreover, $F_j$ is spanning if and only if
$R_j=V_n\setminus T_j$.
\end{lemma}

\begin{proof}
The standard shelling criterion requires, for every $i<j$, an index $r<j$
and a vertex $x\in F_j\setminus F_i$ such that
$F_r\cap F_j=F_j\setminus\{x\}$; see \cite[Definition~2.2]{CS}.
Here $F_j\setminus F_i=T_i\setminus T_j$, while
\[
  F_r\cap F_j=F_j\setminus\{x\}
  \quad\Longleftrightarrow\quad
  T_r=(T_j\setminus\{y\})\cup\{x\}
\]
for some $y\in T_j$.  This gives \eqref{eq:exchange}.  Applying the same
condition to every $x\in F_j$ gives the spanning criterion.
\end{proof}

Let $L_n$ be the order obtained by specializing to $p=3$ the explicit order
of Chauhan and Shukla in Section~3 of \cite{CS}, culminating in their
Definition~3.10.  We use only the defining formulas and verify the resulting
finite lists directly.  This is deliberate: Section~3 of \cite{CS} is written
under the hypothesis $n\ge6p-3$, that is $n\ge15$ at $p=3$, so at $n=13,14$
its formulas are being evaluated outside the range in which \cite{CS} use
them, and none of the lemmas proved there is invoked below.  That the blocks
partition the facet set, that $L_n$ lists every disconnected triple exactly
once, and that the exchange condition holds are all checked on the two finite
lists themselves.  All positions below are one-based.

For $n=13$, replace positions $112$--$117$ of $L_{13}$ by
\begin{equation}\label{eq:r13}
\begin{aligned}
(&\triple{3}{5}{12},\triple{4}{5}{11},\triple{4}{5}{12},
  \triple{5}{9}{12},\triple{5}{10}{12},\triple{5}{11}{12})
\quad\longmapsto\quad{}&\\[-0.3em]
(&\triple{4}{5}{11},\triple{5}{11}{12},\triple{5}{10}{12},
  \triple{5}{9}{12},\triple{3}{5}{12},\triple{4}{5}{12}).&
\end{aligned}
\end{equation}
For $n=14$, replace positions $146$--$155$ of $L_{14}$ by
\begin{equation}\label{eq:r14}
\begin{aligned}
(&\triple{3}{5}{13},\triple{4}{5}{11},\triple{4}{5}{12},
  \triple{4}{5}{13},\triple{5}{9}{12},\triple{5}{9}{13},\\[-0.3em]
 &\triple{5}{10}{12},\triple{5}{10}{13},\triple{5}{11}{12},
  \triple{5}{11}{13})
\quad\longmapsto\quad{}&\\[-0.3em]
(&\triple{4}{5}{11},\triple{4}{5}{12},\triple{5}{9}{12},
  \triple{5}{10}{12},\triple{5}{11}{12},\triple{5}{11}{13},\\[-0.3em]
 &\triple{5}{10}{13},\triple{5}{9}{13},\triple{3}{5}{13},
  \triple{4}{5}{13}).&
\end{aligned}
\end{equation}
Denote the repaired orders by $\widehat L_{13}$ and $\widehat L_{14}$.
Both edits are permutations of a block: the facet set, and the set of triples
occurring before any position outside the block, are unchanged.  Since an
internal permutation of a block changes neither the set of predecessors of a
position outside the block nor the set of earlier triples seen from it,
condition \eqref{eq:exchange} and the spanning criterion can change only at
the $6$ and $10$ edited positions.  In particular no facet outside the two
windows is affected, and the spanning facets listed in
Proposition~\ref{prop:certificates} all lie outside them.

The unrepaired orders fail Lemma~\ref{lem:criterion} only through the
following obstruction pairs:
\[
\begin{array}{c@{\qquad}c@{\qquad}c@{\qquad}l}
\toprule
n & j & T_j & \text{earlier witness triples }T_i\\
\midrule
13 & 112 & \triple{3}{5}{12} & \triple{2}{5}{9}\\
13 & 114 & \triple{4}{5}{12} & \triple{2}{5}{9},\ \triple{2}{5}{10}\\
14 & 146 & \triple{3}{5}{13} & \triple{2}{5}{9}\\
14 & 149 & \triple{4}{5}{13} & \triple{2}{5}{9},\ \triple{2}{5}{10}\\
\bottomrule
\end{array}
\]

\begin{proposition}\label{prop:certificates}
The orders $\widehat L_{13}$ and $\widehat L_{14}$ are shelling orders for
$\Delta_3(C_{13}^3)$ and $\Delta_3(C_{14}^3)$, respectively.  Their spanning
facets are the complements of the omitted triples
\[
  \mathcal S_{13}=\{\triple{4}{10}{12}\}
\]
and
\[
\begin{aligned}
  \mathcal S_{14}=\{&\triple{0}{7}{13},\triple{6}{12}{13},
  \triple{1}{8}{13},\triple{5}{12}{13},
  \triple{2}{9}{13},\\
  &\triple{4}{10}{13},\triple{4}{11}{13},
  \triple{3}{10}{13}\}.
\end{aligned}
\]
In particular, the two orders have one and eight spanning facets.
\end{proposition}

\begin{proof}
The verification described in the section \emph{Exact verification} below
reconstructs $L_n$ from the defining formulas in \cite{CS}, checks that the
blocks used there form a disjoint partition of the facet set, and applies
\eqref{eq:r13}--\eqref{eq:r14}.  It then performs the following exact checks
for each $n\in\{13,14\}$:
\begin{enumerate}
\item the listed triples are exactly the disconnected triples, without
      omission or repetition;
\item condition \eqref{eq:exchange} holds for every pair $i<j$;
\item independently, for each $j$, every inclusion-maximal member of
      $\{F_i\cap F_j:i<j\}$ has cardinality $|F_j|-1$;
\item the spanning triples are exactly those displayed above.
\end{enumerate}
The audit is
\[
\begin{array}{c@{\quad}r@{\quad}r@{\quad}c@{\quad}c@{\quad}c}
\toprule
n & \text{facets} & \text{pairs} & \text{exch. fail.} &
\text{purity fail.} & \text{spanning}\\
\midrule
13 & 169 & 14{,}196 & 0 & 0 & 1\\
14 & 238 & 28{,}203 & 0 & 0 & 8\\
\bottomrule
\end{array}
\]
Lemma~\ref{lem:criterion} now proves the proposition.
\end{proof}

\begin{remark}\label{rem:forced}
The two spanning counts are not independent data.  For a pure shellable
$d$-complex the number of spanning facets equals $h_{d+1}$, and this in turn
equals $|\widetilde\chi|$; the reduced Euler characteristics of
$\Delta_3(C_{13}^3)$ and $\Delta_3(C_{14}^3)$ are $-1$ and $+8$, as one
computes directly from their face numbers, with the sign $(-1)^{n-4}$.  Hence
once shellability is known the counts $1$ and $8$ are forced by the two
complexes themselves, independently of any order, and they act as a
consistency check on the verification rather than as a separate claim.  What
is not forced is which facets are spanning: that depends on the order.
\end{remark}

\begin{proof}[Proof of Theorem~\ref{thm:main}]
Proposition~\ref{prop:certificates} treats $n=13,14$.  Since every facet has
$n-3$ vertices, these complexes have dimension $n-4$.  A pure shellable
$d$-complex with $\beta$ spanning facets is homotopy equivalent to a wedge of
$\beta$ copies of $\Sph^d$ \cite[Theorem~12.3]{Kozlov}; see also
\cite[Theorem~2.3]{CS}.  This gives $\Sph^9$ and
$\bigvee^8\Sph^{10}$ in the two boundary cases.  For $n\ge15$, the result is
\cite[Theorem~1.4]{CS}.  Finally,
\[
  \binom{n-6}{2}-20=\frac{n^2-13n+2}{2},
\]
so the multiplicity agrees with \eqref{eq:conjecture} at $p=3$.
\end{proof}

\section*{Exact verification}

Every computation above is finite and exact, carried out in integer
arithmetic on subsets of $\Z/n\Z$; no floating-point arithmetic is involved.
From the facet criterion of Section~2 one enumerates the $169$ and $238$
disconnected triples.  The order $L_n$ is generated from the defining
formulas of \cite[Section~3]{CS} at $p=3$, the blocks \eqref{eq:r13} and
\eqref{eq:r14} are substituted, and condition \eqref{eq:exchange} is then
evaluated on all $14{,}196$ and $28{,}203$ pairs $i<j$, together with the
purity and spanning tests listed in the proof of
Proposition~\ref{prop:certificates}.  As a control, the same four checks were
applied to the order of \cite[Section~3]{CS} at $n=15,16$, inside the range
$n\ge15$ covered by \cite[Theorem~1.4]{CS}, and returned the expected spanning
counts $16$ and $25$.  No data file, and no copy of the program that produced
the two certificates, accompanies this paper: every
quantity quoted above is either printed here or recomputable from the criteria
of Section~2 together with the definitions of \cite[Section~3]{CS}, and that
program and its transcript are retained in an auxiliary archive,
available on request.  The independent verification program described under
\emph{Independent recomputation} below, and its transcript, are distributed
with this note.  What a reader can and cannot check from the data
printed here is set out next.

\paragraph{Status.}
The orders $\widehat L_{13}$ and $\widehat L_{14}$ are not printed here.
Reconstructing them requires the definitions of \cite[Section~3]{CS}
specialized to $p=3$; given those, the two blocks \eqref{eq:r13} and
\eqref{eq:r14} determine the certificates completely.  Consequently the
following are supported by the computation just described, but are not
checkable from this text alone: the zero failure counts of the audit table;
the assertion that the four pairs tabulated in Section~2 are the
\emph{only} obstructions of $L_{13}$ and $L_{14}$; and the identity---as
opposed to the number---of the spanning facets in $\mathcal S_{13}$ and
$\mathcal S_{14}$.  For the same reason the printed data pin the internal
order of each block only up to permutations with the same in-block discharge
pattern, and do not pin the identity of the witness triples $T_i$.  What
\emph{is} checkable from the printed data is the local part of the argument:
that the six and ten listed triples are omitted triples of $C_{13}^3$ and
$C_{14}^3$; that each new block is a permutation of the old one, so that
by the locality observation of Section~2 nothing outside the two windows can
change; that the four tabulated obstructions are discharged by the
reordering, each through the one vertex of $T_i\setminus T_j$ that can ever
enter $R_j$; that all $15$ pairs internal to the $n=13$ window and $44$ of
the $45$ internal to the $n=14$ window are discharged from in-block
predecessors alone, the exception being the pair at positions $147$ and
$154$, which requires a witness from the unprinted prefix; and, once
shellability is granted, the counts $1$ and $8$, which
Remark~\ref{rem:forced} then forces from the face numbers alone.
Shellability itself is not among these checkable items.

\paragraph{Independent recomputation.}
The computation was re-run independently, by a program written against the
statements printed above, and all $138$ of its checks agreed.  Its reach is
narrower than the paper's in three respects.  It does not rebuild
$\widehat L_{13}$ and $\widehat L_{14}$, which cannot be rebuilt from this
text, so the zero failure counts of the audit table and the identity of
$\mathcal S_{13}$ and $\mathcal S_{14}$ as the spanning sets \emph{of those
two orders} are not re-established by it.  What it does re-establish is the
conclusion together
with the repair mechanism: shelling orders for $\Delta_3(C_{13}^3)$ and
$\Delta_3(C_{14}^3)$, constructed independently and verified on all
$14{,}196$ and $28{,}203$ pairs, pure of dimensions $9$ and $10$, with exactly
$1$ and $8$ spanning facets and realizing $\mathcal S_{13}$ and
$\mathcal S_{14}$; and every assertion of Section~2 that depends only on the
printed blocks.  Its census of explicitly constructed shellings covers
$13\le n\le32$, its reduced Euler characteristics $13\le n\le22$, and its
homology with $\mathbb F_2$ coefficients $13\le n\le20$; the range $n\ge15$ of
Theorem~\ref{thm:main} rests on \cite[Theorem~1.4]{CS} and not on computation.
Odd torsion below the top degree is not excluded by that $\mathbb F_2$
computation; it is excluded by shellability together with
\cite[Theorem~12.3]{Kozlov}.  A recomputation that agrees is evidence about
the computation, not a second proof.

\begin{thebibliography}{9}

\bibitem{BayerEtAl}
M.~Bayer, M.~Denker, M.~Jeli\'c Milutinovi\'c, R.~Rowlands, S.~Sundaram,
and L.~Xue,
\emph{Topology of cut complexes of graphs},
SIAM J. Discrete Math. \textbf{38} (2024), no.~2, 1630--1675.
\href{https://doi.org/10.1137/23M1569034}{doi:10.1137/23M1569034}.

\bibitem{CSV}
P.~Chauhan, S.~Shukla, and K.~Vinayak,
\emph{Shellability of $3$-cut complexes of squared cycle graphs},
J. Homotopy Relat. Struct. \textbf{20} (2025), 163--193.
\href{https://doi.org/10.1007/s40062-025-00365-w}
{doi:10.1007/s40062-025-00365-w}.

\bibitem{CS}
P.~Chauhan and S.~Shukla,
\emph{Shellability of $3$-cut complexes of powers of cycle graphs},
\href{https://arxiv.org/abs/2603.23155}{arXiv:2603.23155v1} [math.CO] (2026).

\bibitem{Kozlov}
D.~Kozlov,
\emph{Combinatorial Algebraic Topology},
Algorithms and Computation in Mathematics, vol.~21,
Springer, Berlin--Heidelberg, 2008.
\href{https://doi.org/10.1007/978-3-540-71962-5}
{doi:10.1007/978-3-540-71962-5}.

\end{thebibliography}

\end{document}